% sections/discussion.tex

\subsection{Reconciling Three Decades of Contradictory Evidence}

The most important contribution of this work is not a new experimental observation but a mechanistic resolution of why two incompatible bodies of evidence about the same biological system have accumulated in parallel for thirty years without either side being proven wrong. Metcalf and the instructive tradition~\cite{metcalf1988} were correct at their scale of measurement: in bulk colony assays with many cells pooled, the variance cancels out, mean-field dynamics prevail, and a sharp cytokine-dependent threshold is the natural result. Hoppe \etal~\cite{hoppe2016} and Olsson \etal~\cite{olsson2016} were equally correct at their scale: at single-cell resolution, the same circuit produces fate heterogeneity that cannot be explained by cytokine dose. Both observations emerge from the same stochastic Petri net model by simply changing the initial molecular count---a parameter that experimentalists implicitly vary between their assays. The model therefore resolves the debate by showing that its framing---instructive \emph{or} stochastic---was incorrect: the system is both, depending on the scale at which it is observed.

The Phase~D and Phase~D-ext GCSF/EPO ratio sweeps add a second resolution layer. EPO-only dose-response experiments that find flat commitment probability are characterising the system in a window where the GCSF/EPO ratio barely changes---a trivial special case of the full two-cytokine control surface. The bifurcation revealed by the ratio sweep is not flat: P(ERY) switches abruptly from $\geq$\,0.92 (48--50/50 erythroid) at GCSF\,=\,2.2 to 0/50 myeloid at GCSF\,=\,2.5, with a threshold at GCSF/EPO\,$\approx$\,2.3 (Table~\ref{tab:commitment_D}). This all-or-nothing cliff is the true bifurcation boundary of the bistable switch, and it depends on the ratio of both signals---not on either cytokine alone. The clinical implication is direct: EPO dose escalation in the absence of GCSF normalisation will not cross the bifurcation threshold; what matters is the relative receptor occupancy of the two competing signals.

This resolution has a broader implication for debates about instruction versus stochasticity in other developmental systems. Wherever a bulk assay and a single-cell assay produce incompatible pictures of the same decision, copy-number-dependent scale dependence is a candidate explanation. The general principle---that stochastic and deterministic behaviours are opposite ends of a continuum parameterised by copy number, not mutually exclusive mechanistic alternatives---is well-established in statistical physics but has not been consistently applied to cell fate debates.

\subsection{The Three-Layer Architecture as a Design Principle}

The decomposition of the commitment cascade into three layers is not merely descriptive---it assigns qualitatively different functional roles to different molecular components whose properties were not previously distinguished (Table~\ref{tab:layers_disc}). The reception layer sets a coarse erythroid bias through receptor stoichiometry; it is the system's compass, pointing the cell toward a fate without specifying it. The decision layer is where the fate is actually chosen, through molecular-count noise at very small copy numbers; it is the coin flip embedded in cell biology. The execution layer converts that choice into a permanent chemical record, quickly and irreversibly.

\begin{table}[ht]
\caption{Three-layer architecture as a design principle. Each layer has distinct EPO-dose sensitivity, noise characteristics, and drug-target implications. ``Bet-hedging'' refers to population diversity maintained through decision-layer stochasticity~\cite{losick2008}; ``chromatin record'' refers to the molecular irreversibility inferred from Pina \etal~\cite{pina2012}.}
\label{tab:layers_disc}
\centering
\small
\begin{tabularx}{\linewidth}{@{}lXlXX@{}}
\toprule
\textbf{Layer} & \textbf{Functional role} & \textbf{EPO-dose sensitivity} & \textbf{Noise / variability} & \textbf{Therapeutic handle} \\
\midrule
Reception & Sets erythroid bias (identity, not probability) & EPO-insensitive; set by receptor stoichiometry & Negligible & Receptor surface residence; EPOR/GCSFR ratio \\
\addlinespace
Decision & Stochastic fate specification & EPO-insensitive & Dominant (bet-hedging) & GATA1 mRNA stability; basal PU.1 transcription rate \\
\addlinespace
Execution & Cytokine-maintained GATA1 dominance & EPO-insensitive once triggered & Negligible & GATA1 kinase; PU.1 E3 ubiquitin ligase \\
\bottomrule
\end{tabularx}
\end{table}

This design principle has a broad implication: because all three layers are EPO-dose-insensitive, the system cannot be finely tuned by adjusting EPO concentration. EPO functions as a binary permissive signal---it enables commitment to proceed but does not control its probability or its depth. Fine-grained control over haematopoietic fate therefore must come from elsewhere: from the copy-number regime of the decision layer (set developmentally), from the receptor stoichiometry of the reception layer (potentially targetable), or from nuclear pH acting on the transcription kinetics that connect all three layers (as demonstrated by Phase C). Importantly, Phase~X demonstrates that the execution layer, without an epigenetic chromatin layer, produces only a cytokine-maintained state: removing EPO reverses the erythroid attractor within approximately 1500~s. This is consistent with the observation of Pina~\etal~\cite{pina2012} that an irreversible molecular record is written before any differentiation marker appears---a record that requires chromatin remodelling not present in the current model, and with the bet-hedging principle of Losick and Desplan~\cite{losick2008}: the system deliberately maintains population diversity through decision-layer stochasticity.

\subsection{Nuclear pH as a Modulator of Erythroid Commitment Probability}

The three-layer architecture identified in Phase B is modulated by a fourth environmental variable: nuclear pH. The pH-Hill mechanism in the model is not a calibration artefact; it is a thermodynamically grounded term encoding the pH sensitivity of the mutual-inhibition kinetics governing GATA1 and PU.1 transcription, via an inhibition constant $K_\text{inh}(\text{pH}) = 0.5 \times 10^{0.5(\text{pH}-7.5)}$~mM that decreases with acidification. A systematic sweep across the physiological pH range (7.0--8.0) at near-bifurcation conditions (GCSF\,=\,1.1~mM, EPO\,=\,0.436--0.446~mM, $N$\,=\,50 per condition, 900 replicates total) reveals that alkaline pH modestly increases erythroid commitment probability: aggregated P(ERY) rises from 2.0\,\% at pH\,=\,7.0 to 3.0\,\% at pH\,=\,7.5 and 4.0\,\% at pH\,=\,8.0, a factor of two across the full physiological range. Standard errors overlap at this sample size, so the effect is directional rather than statistically resolved; the EPO* shift induced by pH variation cannot be quantified from the current data.

The mechanism underlying this trend is consistent with the $K_\text{inh}$ formula, but its direction is counter-intuitive. At alkaline pH, $K_\text{inh}$ is larger, meaning GATA1 must reach higher concentrations to suppress PU.1 with the same efficacy. Despite this weakened mutual-inhibition threshold, P(ERY) is slightly higher at alkaline pH near the bifurcation. A separate high-EPO experiment (EPO\,=\,1.0~mM, no GCSF competition, $N$\,=\,50 per pH; Phase~H) resolves this apparent paradox: once cells commit under conditions where EPO greatly exceeds EPO*, the erythroid attractor is deepest at \emph{acidic} pH---residual nuclear PU.1 is nearly zero at pH\,=\,7.0 (mean 0.000092~mM) and rises monotonically to 0.006~mM at pH\,=\,8.0. The smaller $K_\text{inh}$ at acidic pH means that committed GATA1 suppresses PU.1 more completely; the attractor well is deeper, not shallower. These two regimes---near-bifurcation commitment probability at moderate EPO, and attractor depth at high EPO---are therefore controlled by the same kinetic parameter in opposite directions depending on the cytokine context.

The physiological interpretation of these regime differences bears on the clinical observation that EPO hyporesponsiveness is most pronounced in patients with chronic kidney disease who also present with metabolic acidosis. The present model suggests a specific mechanism: at near-bifurcation EPO concentrations, acidic nuclear pH reduces erythroid commitment probability relative to neutral pH, an effect that EPO dose escalation cannot overcome because it acts through the mutual-inhibition kinetics rather than through receptor occupancy. Conversely, in cells that do commit under acidic conditions, PU.1 silencing is more complete. The clinical hyporesponsiveness may therefore reflect a reduced fraction of progenitors committing (probability effect) rather than a defect in those that do commit (depth effect). This prediction---that acidic pH selectively reduces P(ERY) but not the depth of erythroid gene expression in committed cells---is accessible to CyTOF experiments that jointly measure commitment fraction and erythroid marker intensity across pH-manipulated cultures.

From a drug discovery standpoint, pH modulation is a genuine but modest handle on erythroid commitment: a factor of two in P(ERY) across the full physiological pH range, compared with the orders-of-magnitude effect of the GCSF/EPO ratio identified in Phase D. Interventions that normalise intracellular pH in acidotic bone marrow niches---bicarbonate supplementation, carbonic anhydrase modulators---could in principle increase commitment probability by this modest factor, but would simultaneously reduce attractor depth (more residual PU.1 in committed cells at alkaline pH). The net therapeutic effect is therefore ambiguous and would need to be evaluated against the clinical context.

\subsection{Five Falsifiable Experimental Predictions}

The value of this model for the experimental community lies in generating non-trivial predictions that would not have been proposed without it. We state five predictions in decreasing order of the conceptual novelty they require.

The first prediction is the EPO pulse-duration assay. The Phase~E simulation sweep has already self-tested this prediction within the model across all six originally planned conditions ($t_{\mathrm{pulse}} \in \{60, 120, 300, 500, 1000, 1800\}$~s), yielding a non-trivial non-monotone result: only the 120~s pulse produces stable post-withdrawal erythroid fate (GATA1$_\text{nuc}=2.17$, PU.1$_\text{nuc}=0.31$ at $t=3600$~s; Table~\ref{tab:phaseE}). Both shorter pulses (60~s---insufficient PU.1 erasure) and longer pulses (300--1800~s---GATA1 overshoot decaying past the tipping point) revert to the myeloid attractor. The ghost-attractor reversal timescale scales systematically with pulse length ($\sim$2000~s for 300--500~s pulses; $\sim$2500~s for 1000~s; $\sim$3000~s for 1800~s), a quantitative signature arising from the larger GATA1 overshoot produced by longer stimulation. The model therefore predicts a specific narrow pulse window---not a simple minimum threshold---in which stable priming without ongoing receptor occupancy is achievable. The experimental prediction is now fully bracketed: the lower bound of the stable ERY window lies between 90 and 120~s (the 90~s trajectory is a near-miss --- GATA1$_\text{nuc}$ reaches 2.31~mM at $t=750$~s but PU.1 is already rebuilding and the system reverts at $t\approx$1000~s); the upper bound lies between 150 and 300~s (150~s is a marginal ERY --- GATA1 decays to 1.59~mM at $t=1000$~s before self-activation recovers it to 2.70~mM by $t=3600$~s; 300~s reverts). The confirmed ERY-maintaining interval is $t_{\mathrm{pulse}} \in [120,\,150]$~s. Cells stimulated for $\leq$90~s or $\geq$300~s should revert, with later reversion for longer stimulation; only cells in the (90,\,300)~s range should maintain erythroid gene expression after cytokine withdrawal. This window-plus-reversal-delay prediction is a direct, falsifiable consequence of the bistable architecture and the EPOR-dependent GATA1 self-activation kinetics; no existing model or experimental framework has predicted or observed this pulse-length non-monotonicity. Once a chromatin remodelling layer is incorporated, an analogous assay would identify the pulse length at which the epigenetic record is written, permitting true irreversible commitment independent of cytokine withdrawal.

The second prediction addresses the EPO-insensitivity of the execution layer. Within cells that have committed to the erythroid fate, the kinetics of GATA1 nuclear accumulation after commitment should be identical at EPO = 0.1, 0.5, 1.0, and 2.0~U/mL. Only the fraction of cells committing varies with EPO dose; for those that do commit, the depth and speed of the erythroid programme are EPO-dose-independent. This prediction can be tested by dual-reporter live imaging (GATA1-GFP/PU.1-mCherry knock-in cells) measuring commitment kinetics per cell across EPO doses, sorting for committed cells only, and testing whether GATA1 nuclear accumulation speed and final level are EPO-dose-independent. This is an experiment that no existing model has motivated.

The third prediction is the direct demonstration of the EPO* artefact. A clonal assay in which N=1, N=10, N=50, N=500 cells per well are stimulated across an EPO dose range should show the apparent commitment threshold sharpening with N and disappearing at N=1. The model predicts precisely the shape of this N-dependence: each single cell has a per-cell commitment probability that is a shallow and nearly flat function of EPO dose, so that the standard error of the proportion falls below 5\% only for N $\gtrsim$ 200. This experiment would directly demonstrate that EPO* research has been characterising a measurement scale property rather than a cellular mechanism.

The fourth prediction concerns receptor stoichiometry as a dominant control variable. A titrated EPOR overexpression (doxycycline-inducible, 0.5$\times$ to 4$\times$ endogenous receptor density) at fixed EPO dose should change P(erythroid) more steeply than an equivalent-fold change in EPO concentration. The model quantifies this: doubling EPOR density shifts the EPOR/GCSFR ratio from 4:1 to 8:1, a qualitatively different signal hierarchy, while doubling EPO concentration shifts EPOR occupancy by only a few percent. This prediction would motivate revisiting the interpretation of experiments in AML and MDS subtypes where receptor expression is dysregulated, and it would suggest a rational design principle for erythropoiesis-stimulating agent development: targeting EPOR surface residence time rather than ligand concentration.

The fifth prediction is that phosphorylated GATA1 fraction (pGATA1\_Ser310/total GATA1) measured at 30 minutes post-stimulation predicts individual cell commitment outcome better than total GATA1 protein level. Because pGATA1 is the active form in execution-layer reactions and total GATA1 includes both the phosphorylated and unphosphorylated pools, pGATA1 fraction is more mechanistically proximal to the commitment ratchet. A single-cell CyTOF or imaging mass cytometry experiment measuring both species at 30 minutes and then tracking final fate at 48 hours would provide a direct test of this prediction.

\subsection{Drug Discovery Implications}

The three-layer model has direct implications for therapeutic strategies targeting erythropoiesis. The finding that the execution layer is EPO-dose-insensitive mechanistically explains a clinical observation that has lacked a molecular basis: EPO hyporesponsiveness in myelodysplastic syndrome, chronic kidney disease, and chemotherapy-induced anaemia, where increasing EPO dose beyond a certain point does not proportionally increase reticulocyte production. If P(erythroid) is bounded at 60--70\% by decision-layer noise, and if that noise is set by intracellular molecular copy numbers rather than by EPO concentration, then the ceiling cannot be raised by cytokine dose escalation. Rational therapy for EPO hyporesponsiveness must therefore address the decision layer---stabilising GATA1 mRNA, reducing GATA1 protein turnover, or reducing the basal PU.1 transcription rate---rather than the reception layer.

The execution layer itself (GATA1 kinase and PU.1 E3 ubiquitin ligase) represents a class of drug targets not previously highlighted in this context. These enzymes are activated within 48 seconds and are EPO-dose-insensitive, meaning they function as autonomous commitment amplifiers once engaged. Inhibiting GATA1 phosphorylation (for example, through the kinase that mediates Ser310 modification) would prevent execution-layer engagement and shift the decision outcome even in cells where the decision layer had begun to tip toward erythroid---a distinct mechanism from targeting upstream transcription or receptor signalling. Similarly, PU.1 stabilisation through USP inhibitors (preventing ubiquitin-mediated degradation) would maintain a myeloid-competent state even in the presence of EPOR signalling, an approach relevant to forced myeloid differentiation in AML therapy.

\subsection{Limitations}

Several limitations of the present work must be stated clearly. The model makes quantitative predictions based on calibrated parameters, but the parameter values are not individually measured from the cell line being modelled; they are derived from published estimates across multiple experimental systems and adjusted for thermodynamic consistency. The model should therefore be considered a hypothesis-generating tool rather than a quantitatively precise replica of any specific cell type. In particular, the predicted priming-crossover time of 146~s is a model-specific number that could shift by a factor of two or more with different parameter choices; the qualitative prediction---that there exists an early cytokine-dependent dominance point in the GATA1/PU.1 cascade---is more robust than the specific timing.

Critically, the model does not include a chromatin remodelling layer. The Phase~X cytokine-withdrawal experiment demonstrates that the erythroid attractor established by EPOR signalling is reversible: EPO removal at $t=500$~s produces complete fate reversal (GATA1$_\text{nuc}\approx 0$, PU.1$_\text{nuc}=1.653$) by $t=3600$~s. True transcriptional irreversibility in vivo is mediated by histone methylation, chromatin accessibility remodelling, and locus-specific DNA methylation that deposit a memory of commitment independent of ongoing receptor occupancy. Adding an epigenetic layer is the highest-priority structural extension of the current model.

The commitment threshold used to classify replicates (GATA1\_nuc $>$ PU.1\_nuc at $t = 3600$~s) is phenomenologically motivated by the bistable structure of the model but has not been calibrated against a specific differentation marker. The fraction P(erythroid) computed here corresponds to a model outcome, not to a specific experimental endpoint such as CD71$^+$/GPA$^+$ surface marker expression. Translation from model outcome to measurable marker requires additional calibration against experimental data.

Finally, the model does not yet reproduce the distribution of commitment times observed in live-cell imaging experiments. Hoppe \etal~\cite{hoppe2016} measured the time from GATA1-GFP reporter expression onset to binary fate segregation and observed a distribution spanning several hours. Reproducing this distribution quantitatively would constitute the most stringent validation of the model and is the highest-priority next step.

\subsection{Relationship to Prior Computational Work}

The present model differs from Huang \etal~\cite{huang2007} and Chickarmane \etal~\cite{chickarmane2009} in three key respects: it includes stochastic dynamics rather than deterministic ODEs, it integrates the full signalling cascade from receptor to nuclear phosphorylation rather than abstracting the circuit to a two-ODE toggle, and it identifies the execution layer as structurally distinct from the decision layer by its EPO-dose-insensitivity. The layer decomposition does not appear in the prior literature and constitutes the principal analytical contribution of this work. The assignment of stochastic noise to the decision layer specifically, and not to the execution layer, is also a new claim that is testable experimentally (predictions three and five above).
